\subsection{Spin-flip spectral pairing and the pure-state obstruction}
\label{sec:seven-spinflip-spectral}

We first bound opposite-pair sums in the marginal spectrum, then combine these bounds across overlapping cuts.
For a normalised three-qubit state \(\rho\), write its
eigenvalues as \(\lambda_1\ge\cdots\ge\lambda_8\).
Put
\[
 e(\rho)=8\operatorname{tr}\rho^2-1,\quad
 b(\rho)=\sum_{1\le i<j\le3}\sum_{a,b=1}^3
    \bigl[\operatorname{tr}(\rho\,\sigma_a^{(i)}
                       \sigma_b^{(j)})\bigr]^2,\quad
 v(\rho)=\sum_{i=1}^4
   \left(\lambda_i+\lambda_{9-i}-\frac14\right)^2.
\]
The convention for \(b\) sums precisely the
weight-two Pauli squares on these three sites;
one- and three-body terms do not enter.

\begin{proposition}[Spin-flip spectral pairing; SP7-1]
\label{prop:seven-spinflip-spectral-pair}
Every three-qubit density matrix satisfies the
\emph{sharp spectral inequality}
\begin{equation}\label{eq:seven-spinflip-pair}
 \boxed{\qquad b(\rho)\ge4v(\rho).\qquad}
\end{equation}
Equality holds, in particular, for every
\[
 \rho_r=(1-r)\frac{I_8}{8}+r\,|g\rangle\langle g|,
 \qquad 0\le r\le1,\quad |g\rangle\in(\mathbb C^2)^{\otimes3}
 \text{ pure},
\]
where \(e(\rho_r)=7r^2\),
\(b(\rho_r)=3r^2\), \(v(\rho_r)=3r^2/4\).
For a pure seven-qubit state, letting \(e_S=e(\rho_S)\)
and \(R=\sum_{|T|=3}e_T\), there is the stronger
\emph{joint spectral compatibility} inequality
\begin{equation}\label{eq:seven-spinflip-cross-cut}
 \boxed{\qquad e_S+16v(\rho_S)\le R
       \qquad\text{for all }|S|=3.\qquad}
\end{equation}
The latter also holds for mixed global states.
\end{proposition}
\begin{proof}
Let \(\widetilde\rho=Y^{\otimes3}\rho^T Y^{\otimes3}\)
be the spin-flipped marginal. This is positive,
has the same eigenvalues as \(\rho\), and flips the
sign of each odd Pauli weight. Pauli orthogonality gives
\[
 \operatorname{tr}\rho^2=\frac{1+e}{8},\qquad
 \operatorname{tr}(\rho\widetilde\rho)
      =\frac{1-e+2b}{8},
\]
hence
\(\operatorname{tr}\rho^2+
  \operatorname{tr}(\rho\widetilde\rho)=(1+b)/4\).
The von Neumann trace rearrangement inequality implies
\[
 \operatorname{tr}(\rho\widetilde\rho)
   \ge\sum_{i=1}^8\lambda_i\lambda_{9-i}.
\]
Consequently
\[
 \frac{1+b}{4}\ge
 \sum_{i=1}^4(\lambda_i+\lambda_{9-i})^2
 =\frac14+v(\rho),
\]
proving~\eqref{eq:seven-spinflip-pair}.

For every pure three-qubit vector \(g\), complementarity
of one- and two-qubit marginal purities gives
\(\sum_{i<j}\operatorname{tr}\rho_{ij}(g)^2
 =\sum_i\operatorname{tr}\rho_i(g)^2\).
Its Pauli expansion shows that the total weight-two
Pauli sum is exactly \(3\), independently of \(g\).
Depolarisation by \(r\) multiplies all nontrivial
Pauli expectations by \(r\), proving \(b=3r^2\)
and \(e=7r^2\). The spectrum is
\(((1+7r)/8,((1-r)/8)^7)\), whose four opposite
pair sums give \(v=3r^2/4\); equality follows.

For a seven-qubit global state, let
\(b_{ij}\) denote the sum of weight-two Pauli
squares on the physical pair \(ij\). Then
\(b(\rho_S)=\sum_{\{i,j\}\subset S}b_{ij}\)
and each such pair occurs in \emph{four other}
three-site subsets. Since all weights in every
\(e_T\) are nonnegative,
\[
 R-e_S\ge4b(\rho_S)\ge16v(\rho_S).
\]
This proves~\eqref{eq:seven-spinflip-cross-cut}.
\end{proof}

\begin{lemma}[Sharp entropy--purity slope; proved, SP7-5]\label{lem:seven-entropy-linear-purity}
Every $d$-dimensional density matrix, $d\ge3$, satisfies
\begin{equation}\label{eq:entropy-purity-general-slope}
 \log_2d-S(\rho)\ge c_d\bigl(d\operatorname{tr}\rho^2-1\bigr),
 \qquad c_d=\frac{(d-1)\log_2(d-1)}{d(d-2)}.
\end{equation}
The constant is sharp: equality holds exactly at the uniform spectrum and at
$\bigl((d-1)/d,(1/[d(d-1)])^{d-1}\bigr)$. In particular, for $d=8$,
\begin{equation}\label{eq:seven-entropy-linear-purity}
 3-S(\rho)\ge\frac{7\log_27}{48}(8\operatorname{tr}\rho^2-1)
 \ge\frac{49}{120}(8\operatorname{tr}\rho^2-1)
 \ge\frac25(8\operatorname{tr}\rho^2-1),
\end{equation}
and the last two comparisons are strict unless $\rho=I_8/8$.
These bounds include rank-deficient matrices.
\end{lemma}
\begin{proof}
Put $m=d-1\ge2$. By Lemma~\ref{lem:seven-fixed-purity-envelope} it suffices to treat the one-spike spectrum
$((1+mr)/d,((1-r)/d)^m)$, $0\le r\le1$. Its normalised purity excess is $d\operatorname{tr}\rho^2-1=mr^2$ and its
entropy loss is $\log_2d-S=F_m(r)/(d\ln2)$, where
\[
 F_m(r)=(1+mr)\ln(1+mr)+m(1-r)\ln(1-r).
\]
The derivative of $F_m(r)/r^2$ has numerator $N=rF_m'-2F_m$, with $N(0)=N'(0)=0$ and
\[
 N''=rF_m'''=\frac{m(m+1)r\bigl(2mr-(m-1)\bigr)}{(1-r)^2(1+mr)^2}.
\]
So $N''$ changes sign once, from negative to positive: $N'$ is negative and then increasing, and $N$ has at most one
interior zero. At $r_*=(m-1)/m$ one has $1+mr_*=m$ and $1-r_*=1/m$, so
\[
 F_m(r_*)=(m-1)\ln m,\qquad F_m'(r_*)=2m\ln m,\qquad N(r_*)=0.
\]
Hence $F_m/r^2$ strictly decreases on $(0,r_*]$ and strictly increases on $[r_*,1]$, with minimum $m^2\ln m/(m-1)$.
Dividing by $dm\ln2$ gives~\eqref{eq:entropy-purity-general-slope}, with equality only at $r\in\{0,r_*\}$.
For $d=8$, $c_8=7\log_27/48$; since $7^5=16807>16384=2^{14}$, $\log_27>14/5$ and
$7\log_27/48>49/120>2/5$. SP7-5 checks these identities and constants.
\end{proof}

\begin{corollary}[Sharpened universal high-entropy spectral
restriction; proved, SC-2, SP7-6]
\label{cor:seven-entropy-eigen-cap-497}
For any (pure or mixed) seven-qubit state \(\omega\),
\[
 E_3(\omega)\ge\frac{59307}{20000}
 \quad\Longrightarrow\quad
 \boxed{\ \lambda_{\max}(\omega_S)
      <\frac{497}{1000}=0.497
      \quad\text{for every }|S|=3.\ }
\]
Every competitor satisfies the hypothesis, since
\(E_*\ge E_3(\psi_{7/100})>59307/20000\) (SC-2).
\end{corollary}
\begin{proof}
Fix a cut \(S\) with eigenvalues in descending
order. Let \(t=\lambda_1\), \(a=\lambda_8\)
and \(b=(1-t-a)/6\), so \(0\le a\le(1-t)/7\).
By concavity of entropy,
\[
 S(\omega_S)\le
 -t\log_2t-a\log_2a-6b\log_2b.
\]
For the four opposite-pair sums of eigenvalues
in \(v(\omega_S)\), one sum is \(t+a\), and
the other three sum to \(1-t-a\).
Cauchy--Schwarz thus yields
\[
 v(\omega_S)\ge\frac43
      \left(t+a-\frac14\right)^2.
\]
The spin-flip compatibility inequality
\eqref{eq:seven-spinflip-cross-cut} and
Lemma~\ref{lem:seven-entropy-linear-purity}
(with the rational slope \(49/120\))
applied to the other 34 marginals yield
\[
\begin{split}
 35(3-E_3(\omega))
 &\ge 3-S(\omega_S)
      +\frac{49}{120}\sum_{T\ne S}e_T\\
 &\ge 3-S(\omega_S)+\frac{98}{15}v(\omega_S)
 \ge J(t,a),
\end{split}
\]
where, continuously at \(a=0\),
\[
 J(t,a)=3+t\log_2t+a\log_2a+6b\log_2b
      +\frac{392}{45}
           \left(t+a-\frac14\right)^2.
\]
Apply the tangent inequality \(p\ln(p/q)\ge p-q\)
(\(p\ge0,q>0\)) to the three probability weights
\((t,6b,a)\) with comparison weights
\(q_1=1/2\), \(6q_2=1/2\) (\(q_2=1/12,\;8q_2=2/3\)), and
\(q_3=1/192\) (\(8q_3=1/24\)), whose sum is
\(q_1+6q_2+q_3=1+1/192\). Since
\(3+t\log_2t+a\log_2a+6b\log_2b
 =\bigl[t\ln(8t)+6b\ln(8b)+a\ln(8a)\bigr]/\ln2\)
and \(\ln(8q_1)=\ln4=2\ln2\),
\(\ln(8q_2)=\ln(2/3)=\ln2-\ln3\),
\(\ln(8q_3)=\ln(1/24)=-3\ln2-\ln3\), this gives
\[
\begin{aligned}
 3+t\log_2t+a\log_2a+6b\log_2b
 &\ge 2t+(1-t-a)(1-\log_23)-a(3+\log_23)
      -\frac1{192\ln2}\\
 &=1+t-(1-t)\log_23-4a-\frac1{192\ln2}.
\end{aligned}
\]
Completing the square in \(a\) gives
\[
 -4a+\frac{392}{45}\left(t+a-\frac14\right)^2
 =4t-1-\frac{45}{98}
  +\frac{392}{45}\left(t+a-\frac{47}{98}\right)^2
 \ge 4t-1-\frac{45}{98}.
\]
The resulting lower bound \(5t-\frac{45}{98}-(1-t)\log_23-\frac1{192\ln2}\) increases in \(t\)
(slope \(5+\log_23>0\)). For \(t\ge t_0=497/1000\), using only
\(3^5=243<256=2^8\) (\(\log_23<8/5\)) and
\(\ln2=2\operatorname{arctanh}(1/3)>2/3\)
(\(1/(192\ln2)<1/128\)), we obtain
\[
 J(t,a)
 >5t_0-\frac{45}{98}-\frac85(1-t_0)-\frac1{128}
 =\frac{4755759}{3920000}
 >\frac{4851}{4000}
 =35\left(3-\frac{59307}{20000}\right),
\]
with exact rational margin \(1779/3920000>0\) (SP7-6),
contradicting \(E_3(\omega)\ge59307/20000\).
Therefore \(t<497/1000\) on all 35 cuts.
\end{proof}

